\chapter{High Temperature}

\section*{Definition and Overview}
A high temperature, or fever, is a body temperature above the normal range, generally taken as 38 degrees Celsius or higher. Fever is not an illness in itself but a sign that the body's immune system is responding to an infection or another trigger. Most fevers are caused by viral infections, are short-lived, and resolve on their own. The height of the fever does not always reflect the severity of the illness, and how unwell a person feels matters more than the number on the thermometer. Treatment aims at comfort, and most fevers in otherwise healthy people need no specific medicine.

\section*{Epidemiology}
Fever is one of the most common reasons for medical consultation and for parents to seek advice about children. Most children have several febrile illnesses in their early years, usually from viral infections such as colds, flu, and gastroenteritis. Fever is more concerning in young infants, in whom serious bacterial infection is more likely. In adults, fever is commonly due to respiratory, urinary, or gastrointestinal infections, while persistent or recurring fever warrants investigation for other causes.

\section*{Aetiology and Pathophysiology}
\begin{itemize}
    \item \textbf{Infection:} the commonest cause; viruses and bacteria trigger immune cells to release pyrogens
    \item \textbf{Pyrogens:} cytokines such as interleukins act on the brain's thermostat in the hypothalamus, raising the set point
    \item \textbf{Immunisation:} vaccines can cause mild, brief fever
    \item \textbf{Inflammation:} autoimmune disease, inflammatory conditions, and tissue injury
    \item \textbf{Other causes:} certain cancers, blood clots, and drugs
    \item \textbf{Heat-related:} environmental heat causes hyperthermia, which is different from fever
    \item \textbf{Pathophysiology:} raising body temperature enhances immune function and inhibits some pathogens; the body conserves and produces heat until the set point is reached
\end{itemize}

\section*{Clinical Features}
\begin{itemize}
    \item Temperature of 38 degrees Celsius or higher
    \item Feeling hot, flushed, shivering, and sweating as temperature rises and falls
    \item Headache, muscle aches, fatigue, and poor appetite
    \item Associated symptoms point to the cause: cough, sore throat, vomiting, diarrhoea, or urinary symptoms
    \item Children may be irritable, clingy, or drowsy
    \item Fever without other symptoms is common, especially in children
    \item Red flags: stiff neck, rash that does not fade, difficulty breathing, seizures, or drowsiness
\end{itemize}

\section*{Differential Diagnosis}
\begin{itemize}
    \item Viral infections are the most common cause
    \item Bacterial infections: chest, urinary, ear, throat, and bloodstream infections
    \item Young infants with fever need urgent assessment for serious infection
    \item Non-infectious causes: inflammatory disease, medications, and malignancy
    \item Hyperthermia from heat exposure, which is not fever
    \item The pattern, duration, and associated symptoms guide the investigation
\end{itemize}

\section*{When to Seek Medical Care}
See a doctor for fever in babies under 3 months, fever lasting more than 3 days, fever with breathing difficulty, rash, severe pain, or dehydration, or if the person is very unwell despite treatment.

\textbf{Call 000 (or local emergency services) for:}
\begin{itemize}
    \item A baby under 3 months with a temperature of 38 degrees or higher
    \item Fever with a rash that does not fade when pressed
    \item Stiff neck, severe headache, or sensitivity to light
    \item Difficulty breathing, rapid breathing, or blue lips
    \item Seizures, confusion, or drowsiness
    \item Fever with an inability to drink or signs of severe dehydration
\end{itemize}

\section*{Diagnosis and Investigations}
\begin{itemize}
    \item \textbf{Temperature measurement:} accurate measurement, using a thermometer rather than touch
    \item \textbf{History and examination:} localising symptoms and examination of ears, throat, chest, abdomen, and skin
    \item \textbf{Urine testing:} common in children with unexplained fever
    \item \textbf{Blood tests:} full blood count, inflammatory markers, and cultures when serious infection is suspected
    \item \textbf{Chest X-ray and other imaging:} as indicated by symptoms
    \item \textbf{Assessment of babies:} specialist evaluation for young infants
\end{itemize}

\section*{Management}
\begin{itemize}
    \item \textbf{Comfort measures:} light clothing, fluids, and rest
    \item \textbf{Paracetamol or ibuprofen:} for discomfort, not routinely for the number on the thermometer
    \item \textbf{Avoid aspirin in children:} it can cause Reye syndrome
    \item \textbf{No sponging with cold water:} it is uncomfortable and ineffective
    \item \textbf{Treat the cause:} antibiotics for bacterial infections, supportive care for viruses
    \item \textbf{Watch and reassess:} seek medical care if symptoms worsen or do not settle
    \item \textbf{Hospital care:} for infants, severe dehydration, or serious infection
\end{itemize}

\section*{Complications}
\begin{itemize}
    \item Febrile seizures in children aged 6 months to 5 years, which are usually brief and harmless
    \item Dehydration from poor intake and fluid loss
    \item Serious bacterial infection may be masked by fever in young infants
    \item Prolonged fever can cause weight loss and fatigue
    \item Delayed diagnosis of treatable infections
    \item Over-treatment with medicines, with risk of side effects
\end{itemize}

\section*{Prognosis and Outcomes}
Most fevers resolve within a few days as the underlying infection clears, and treatment is directed at comfort rather than the temperature itself. Serious causes are uncommon, especially in older children and adults, and are identified by the pattern and associated symptoms. With appropriate fluid intake and monitoring, almost all people recover fully without complications. The key is assessing how unwell the person is, rather than focusing on the height of the fever.

\section*{Prevention and Self-Care}
\begin{itemize}
    \item Keep a working thermometer and measure temperature accurately
    \item Offer plenty of fluids to prevent dehydration
    \item Dress lightly and keep the room comfortable
    \item Use paracetamol or ibuprofen for discomfort according to age and weight
    \item Monitor for red flags: drowsiness, rash, breathing difficulty, and poor drinking
    \item Seek urgent care for fever in babies under 3 months
    \item Keep vaccinations up to date to prevent febrile illnesses
    \item Practise good hand hygiene to reduce infections
\end{itemize}

\section*{Sources and Further Reading}
The following reputable sources provide further information for healthcare students and patients seeking reliable, current advice about high temperature.

\begin{itemize}
    \item Healthdirect. \textit{Fever: what it means and what to do.} https://www.healthdirect.gov.au/fever
    \item Raising Children Network. \textit{Fever in children.} https://raisingchildren.net.au/
    \item Royal Children's Hospital Melbourne. \textit{Fever fact sheet.} https://www.rch.org.au/
    \item NHS. \textit{High temperature (fever) in adults and children.} https://www.nhs.uk/conditions/fever/
    \item Mayo Clinic. \textit{Fever: first aid.} https://www.mayoclinic.org/
    \item MedlinePlus. \textit{Fever.} https://medlineplus.gov/fever.html
\end{itemize}
